\section{Record excerpts}
\label{app:records}

Real, trimmed excerpts from the committed artifacts, checked against the files named below on
2026-09-29: a claim record, the N1 freeze file, and the two GDPR traces moved here from
\cref{sec:demos}. Each is what the code wrote, not a reconstruction for this report; long fields are
shortened with \ldots{} where noted. Full discussion: \cref{sec:evidence-model,sec:n3}.

\subsection{A claim record, with its evidence selector}

\repo{reconstruct/runs/20260928T145943Z-bc1cc9a30a02/ledger.json} (GDPR v2), one claim, trimmed:

{\small
\begin{verbatim}
{
  "assertion": "A DPIA should attach any relevant additional
    documents referenced in the DPIA, such as Privacy Notices
    and consent documents.", "knowledge_type": "check",
  "evidence": [{
    "selector": {
      "exact": "attached any relevant additional documents we
        reference in our DPIA, e.g. Privacy Notices, consent
        documents",
      "locator": "sha256:14b9bb35...f8a23e7;char=9061,9169"
    },
    "source": {"identifier": "https://ico.org.uk/.../dpias/",
        "kind": "procedure_document", "world": "public"},
    "verification": {"verdict": "pending", "verifier": null}
  }],
  "generation": {"activity": "extraction", "agent": {
    "family": "anthropic", "id": "anthropic/claude-haiku-4.5"}}
}
\end{verbatim}
}

The \code{assertion} is the claim in the extractor's own words; \code{evidence[].selector} ties it
to a verbatim span (\code{exact}) and to a character range in a hashed, cached document
(\code{locator}), so the span can be re-sliced later without trusting the extractor.
\code{verification.verdict} is \code{pending} because this is one of the 140 GDPR v2 evidence items
whose verifier call failed silently. The claim's computed label is therefore
\code{synthetic_extrapolation}: only a \code{supports} verdict from the second model family makes a
claim \code{literature_supported} (\cref{sec:evidence-model}).

\subsection{The N1 freeze file}

\repo{residual/frozen/n1.json}, in full:

{\small
\begin{verbatim}
{
  "schema_sha256":
    "f2ffcfdfab6ea55a2e187c864d14a411e1a6cd2b2c54bd0fa7b9dbfb70652591",
  "features_sha256":
    "f2d3e07999560b5587987d31e96353a4fd526acf600dabaebc4d79cb54a3bfe9",
  "semantics_sha256":
    "f3e85218b9d3868467ed7c6ca646496ccfebcd6d91943ce152c7884b3b391d5d",
  "frozen_by": "N1"
}
\end{verbatim}
}

Three SHA-256 hashes freeze three things: the \code{ClaimRecord} schema, the feature set (the
tuple of \code{AreaFeatures} field names), and the semantics (the code that decides an epistemic
label from evidence). A test (\code{test_the_schema_matches_its_freeze}) fails if the current code's
hash of any of the three no longer matches this file, which forces a deliberate, committed re-freeze
(\cref{app:claims}, row E07).

\subsection{Trace 2: GDPR v2, the boundary of \enquote{relevant controller}}

\textbf{GDPR v2 is not N3-admissible} (140 evidence items never verified); this trace illustrates the
record format only.\footnote{\repo{research/n3/gdpr_v2/gapmap.json}, \code{map[0]}, \code{gap_id}
\code{g-656818235091}.}

\begin{table}[htbp]
  \centering
  \small
  \caption{Trace 2: GDPR v2 (not N3-admissible), lens DISC, gap \code{g-656818235091}, category
  HYP, closure \code{partial}.}
  \label{tab:trace2}
  \begin{tabularx}{\linewidth}{@{}p{2.4cm}X@{}}
    \toprule
    Stage & Content \\
    \midrule
    Sources & Seed: \url{legislation.gov.uk/eur/2016/679/article/35} (two URLs, one independence
      key, statute HTML). Topic: \url{edpb.europa.eu} (Irish DPC DPIA decision, PDF) and
      \url{cnil.fr} (WP248, PDF). Three independent keys discuss the topic. \\
    Seed span (verbatim) & \enquote{Compliance with approved codes of conduct referred to in
      Article 40 by the relevant controllers or processors shall be taken into due account in
      assessing the impact of the processing operations performed by such controllers or
      processors, in particular for the purposes of a data protection impact assessment.} \\
    Seed claims / label & \code{c-45c0ec0f142fd0ba} and \code{c-c08f1d894cabec8e}, both
      \code{literature_supported}, both from the one legislation key \\
    Lens \& target & DISC: the term \enquote{relevant controller} is used, but no boundary is
      stated \\
    Inferred gap & Of 8 verified sentences retrieved, the judge found only a partial statement of a
      boundary: \enquote{Does not explicitly state threshold, defining feature, or contrast.} \\
    Hypothesis (inferred) & \enquote{Hypothesis (inferred): practitioners are predicted to
      discriminate \enquote{relevant controller} from its neighbours by features that no verified
      claim in this ledger matched test DISC:relevant controller for.} \\
    Question \& channel & \enquote{One source writes: [the seed span]. Think of a recent case where
      you had to judge whether something was relevant controller. Describe one case that clearly
      was, one that clearly was not, and one that was hard to call. What differed between them, and
      what did you check first?} -- contrasting-case classification, or expert-rated written
      vignettes \\
    \bottomrule
  \end{tabularx}
\end{table}

\textbf{Weaknesses.} The run is not admissible. The anchor is a legal drafting phrase, not a
judgment term, so this is arguably a DISC false positive. The question is ungrammatical
(\enquote{whether something was relevant controller}), and the hypothesis template ends in a garbled
clause. The state is \code{partial}: the judge found part of a boundary statement. The sources are
statute and regulator text, not practitioner text.

\subsection{Trace 3: GDPR v1's only gap candidate}

GDPR v1's final map holds one gap candidate.\footnote{\repo{research/n3/gdpr_v1/gapmap.json},
\code{map[0]}, \code{gap_id} \code{g-1b1e68dc5d39}.}

\begin{table}[htbp]
  \centering
  \small
  \caption{Trace 3: GDPR v1, lens RESULT, gap \code{g-1b1e68dc5d39}, category HYP, closure
  \code{partial}.}
  \label{tab:trace3}
  \begin{tabularx}{\linewidth}{@{}p{2.4cm}X@{}}
    \toprule
    Stage & Content \\
    \midrule
    Sources & Seed: \url{ico.org.uk} (UK). Topic: \url{dataprotection.ie} (Ireland),
      \url{cnpd.public.lu} (Luxembourg) and further ICO pages. Two independence keys: ICO and CNPD
      share one key under the old 25-word union rule, since dropped. \\
    Seed span (verbatim) & \enquote{The nature of the processing is what you plan to do with the
      personal data. This should include, for example: how you collect the data; how you store the
      data; how you use the data; who has access to\ldots} \\
    Seed claim / label & \code{c-6370fb5b7eeb5c13}, \code{literature_supported}. Assertion:
      \enquote{The nature of processing in a DPIA should include details about data collection,
      storage, usage, access, sharing, processors, retention periods, security measures, new
      technologies, and novel processing types.} \\
    Lens \& target & RESULT, anchored on the assertion's tail \enquote{measures, new technologies,
      and novel processing types.}: a test is prescribed, but not how a practitioner reads its
      result \\
    Inferred gap & Of 8 verified sentences retrieved, the judge found only a partial statement:
      \enquote{Sentence 8 mentions mitigating measures but does not explicitly state expected results
      or next actions.} The lexical rule reads \code{closed}. \\
    Question \& channel & \enquote{Sources prescribe: \enquote{The nature of processing in a DPIA
      should include details about data collection, storage, usage, access, sharing, processors,
      retention periods, security measures, new technologies, and novel processing types.}. Think of
      the last time you did this. What result did you get, what had you expected, what did that
      result make you do next -- and what result would have sent you down a different path?} -- CDM
      probes on a recalled case; process tracing \\
    \bottomrule
  \end{tabularx}
\end{table}

\textbf{Weaknesses.} The text after \enquote{Sources prescribe} is the extractor's assertion, not a
source quote: the code fell back to the paraphrase and still set it in quotation marks. The anchor
is a fragment of the same paraphrase. Only two independence keys stand behind the candidate, one of
them an over-merge of two national regulators. Lexical closure reads \code{closed} where the judge
reads \code{partial}, a small instance of the disagreement reported in \cref{sec:absence-check}.

The record as stored, trimmed:

{\small
\begin{verbatim}
{
  "gap_id": "g-1b1e68dc5d39", "lens": "RESULT", "category": "HYP",
  "missing": "a result-to-interpretation mapping for 'measures,
    new technologies, and novel processing types.'",
  "confidence": {"A": 1, "B": 0, "P": 1, "Q": 0,
    "breadth": "narrow", "k_step": 1, "k_topic": 2,
    "robustness": "0/4", "score": 0},
  "inferred_gap": {"closure_state": "partial",
    "judge_model": "qwen2.5:7b-instruct",
    "test_id": "RESULT:c-6370fb5b7eeb5c13"},
  "hypothesis": {"label": "inferred",
    "predicted_knowledge_type": ["interpretation", "expectancy"],
    "text": "Hypothesis (inferred): practitioners are predicted
      to read the result of 'measures, new technologies, and
      novel processing types.' against expected values and map
      it to the next action; no verified claim in this ledger
      matched test RESULT:c-6370fb5b7eeb5c13."}
}
\end{verbatim}
}

\code{confidence} is the breadth rubric ($A$ = topic breadth, $B$ = step breadth, $P$ =
partial-state penalty, $Q$ = promotional penalty), never a calibrated probability: no gold exists to
calibrate against. \code{hypothesis.label} is always \code{inferred}.
